% !TITLE 浣溪沙·红日已高三丈透
% !AUTHOR 李煜
% !DYNASTY 五代十国
% !GENRE 词
% !ORDER 1
% !EDITION 《南唐二主词笺注》，王仲闻校订、陈书良与刘娟笺注，中华书局（2013年版为候选）
% !EDITION_STATUS pending

\begin{poem}
红日已高三丈透\textsuperscript{1}，
金炉\textsuperscript{2}次第\textsuperscript{3}添香兽\textsuperscript{4}。
红锦地衣\textsuperscript{5}随步皱。

佳人舞点\textsuperscript{6}金钗溜\textsuperscript{7}，
酒恶\textsuperscript{8}时拈\textsuperscript{9}花蕊嗅。
别殿\textsuperscript{10}遥闻箫鼓奏。
\end{poem}

\begin{annotations}
\notetext{1}{三丈透：太阳升起三丈高了。透，完全升起。说明早已日上三竿，宴会从昨夜通宵持续到现在。}
\notetext{2}{金炉：铜制的香炉。宫廷用品，显示奢华。}
\notetext{3}{次第：依次、一个接一个。香炉里的香料烧完了，宫女们依次添加。宴会之久，可见一斑。}
\notetext{4}{香兽：做成兽形的香料（如香饼、香丸），或用兽形铜炉焚烧的香料。唐五代宫廷宴会中常见。}
\notetext{5}{红锦地衣：红色锦缎铺成的地毯。“随步皱”——舞步踏过，地毯起了皱褶。写出舞蹈的投入。}
\notetext{6}{舞点：舞蹈的节拍。随着鼓点起舞，动作越来越快。}
\notetext{7}{金钗溜：头上的金钗滑落下来。舞蹈太尽兴了，连首饰都顾不上。写出沉醉忘我的状态。}
\notetext{8}{酒恶：醉了、酒意上头。“恶”是当时的口语，意思是酒喝多了不舒服。李煜用作“沉醉”的同义词。}
\notetext{9}{拈：用手指轻轻地拿、捏。拈（niān）。微醉时拈起花瓣来闻，动作慵懒娇柔。}
\notetext{10}{别殿：其他宫殿。远处别的宫殿也传来箫鼓声——整座宫殿都在通宵欢宴。这是亡国前最后的繁华。}
\end{annotations}

\begin{appreciation}
五代这一卷里，李煜的词按顺序读，是先甜后苦。这一首写在他还做皇帝的时候，通宵宴饮，天亮未歇，是他早期“富贵”一路的代表。

“红日已高三丈透”，太阳升起三丈高了。不说“日上三竿”，偏要量一量那三丈，是嫌宫中的钟漏走得太慢，一晌贪欢，连天光也不放在眼里。“金炉次第添香兽”，香兽是做成兽形的香，烧尽一块便续上一块，宫人依次来添；香添了几回，便是几段光阴过去了。“红锦地衣随步皱”，猩红的地毯被舞步踩出褶皱，足下一层细细的皱，把满堂的热闹都收在了毯子面上。

下片镜头移向人。“佳人舞点金钗溜”，舞点是鼓点，舞女踩着拍子旋转，鬓边的金钗滑落下来也顾不上扶，可见舞到了忘形。“酒恶时拈花蕊嗅”，酒恶是当时的口语，意思是酒喝过了头，人有点发晕；她随手拈一朵花来闻，想借花香醒一醒酒。一个“拈”字把醉后慵懒娇软的神气全写活了。收句忽然荡开：“别殿遥闻箫鼓奏”，远处别的宫殿里也隐隐传来箫鼓之声。原来彻夜不歇的不止这一处，整座皇宫都在狂欢。

通篇没有一个“我”字，李煜始终不出场，镜头只对着日影、香兽、地毯与金钗，像一台摄影机架在宴席一角。可是越把热闹拍到极处，人心里越发慌：盛到顶点，下一步就是散场。

《生年不满百》里有两句：“昼短苦夜长，何不秉烛游。”白日苦短，那就点起蜡烛把夜晚接着过。李煜几乎正是把这两句诗过成了日子，帐里焚的是香，毯上踏的是舞，一宴接连着一宴。那位无名诗人的秉烛游是乱世里对短促生命的认命，李煜的秉烛游却是太平年头里不知收敛的挥霍。

另有一位北宋词人晏殊，也写过一首《浣溪沙》。同一个词牌，同一个曲调，晏殊填进去的是“无可奈何花落去”，一句一叹息；李煜填进去的却是“红日已高三丈透”，一句一沉醉。词牌只管定曲调，填进去的却是两样的心事。等到许多年后他被押到汴京，在囚室里再回想这样的夜晚，才明白那一晌的欢是用整个江山换来的。
\end{appreciation}
